\section{Density-functional perturbation theory in practice}\label{app:E.4}

A converged run of the self-consistent Sternheimer equation of \cref{sec:5.4} returns the screened potential, from which the vertex follows; the dynamical matrix; and, at $\vect{Q} = 0$, the response to a uniform field.

\subsection{Outputs of a linear-response run}

One perturbation is labeled by a phonon wavevector $\vect{Q}$, a nucleus $\kappa$ and a Cartesian direction $\alpha$. For each perturbation, the Sternheimer equation \cref{eq:5.4-sth}, the first-order density \cref{eq:5.4-dn}~\cite[Eq.~(49)]{giustino2017electron} and the screened potential \cref{eq:5.4-dV} are iterated to self-consistency, exactly as a ground-state calculation is, and the cost of one perturbation is of the same order as the cost of the ground-state calculation~\cite[Sect.~II.C]{baroni2001phonons}.

The screened potential $\partial_{\kappa\alpha,\vect{Q}}V_{\mathrm{KS}}$ is obtained for all $3N_{\mathrm{at}}$ perturbations at this $\vect{Q}$, with $N_{\mathrm{at}}$ the number of atoms per cell, and its matrix elements between Bloch states, weighted by the polarization vectors, are the electron--phonon vertex $g_{mn\nu}(\vect{k},\vect{Q})$ of \cref{eq:A.6-g}. The dynamical matrix $D^{\mathrm{dyn}}(\vect{Q})$ of \cref{eq:A.6-Ddyn} follows from the first-order density by \cref{eq:5.4-hess}. Diagonalizing it yields $\omega_{\vect{Q}\nu}$ and $\vect{e}_{\kappa\nu}(\vect{Q})$. The response to a uniform electric field gives the electronic dielectric tensor $\varepsilon_\infty$ and the Born effective charges $Z^*_\kappa$. DFPT frequencies inherit the strengths and the failures of the ground-state calculation (\cref{app:E.1}).

\subsection{Born charges and dielectric tensors}

A uniform electric field $\vect{\mathcal{E}}$ cannot be written as a periodic potential, since $-e\vect{\mathcal{E}}\cdot\vect{r}$ grows without bound. DFPT avoids the problem because \cref{eq:5.4-sth} needs only matrix elements of the perturbation between different eigenstates, and the matrix elements of $\vect{r}$ between different eigenstates are well defined through the commutator,
\begin{equation}
  \langle\psi_{c\vect{k}}|\vect{r}|\psi_{v\vect{k}}\rangle = \frac{\langle\psi_{c\vect{k}}|[H_{\mathrm{KS}},\vect{r}]|\psi_{v\vect{k}}\rangle}{\varepsilon_{c\vect{k}} - \varepsilon_{v\vect{k}}} ,
  \label{eq:E.4-rcomm}
\end{equation}
which is the relation between the position and velocity matrix elements of \cref{sec:3.2}~\cite[Eq.~(37)]{baroni2001phonons}. For a local potential $[H_{\mathrm{KS}},\vect{r}] = -\ii\vect{p}/m_e$; a nonlocal pseudopotential (\cref{sec:3.5,app:E.2}) adds its own commutator, and omitting it biases every quantity built from the optical limit, the dielectric constant first among them. The projected state $\op{P}^{\mathrm{c}}\vect{r}|\psi_{v\vect{k}}\rangle$ is itself the solution of a Sternheimer equation with $[H_{\mathrm{KS}},\vect{r}]|\psi_{v\vect{k}}\rangle$ on the right-hand side~\cite[Eq.~(44)]{baroni2001phonons}. With the projected state, the field enters \cref{eq:5.4-sth} like any other perturbation, and the self-consistent solution for the field perturbation is the induced macroscopic polarization $\vect{P}$.

The macroscopic quantities then follow by differentiation. In the Heaviside--Lorentz units of the book (\cref{app:A.2}), with $\vect{D} = \vect{\mathcal{E}} + \vect{P}$,
\begin{equation}
  \varepsilon_\infty^{\alpha\beta} = \delta_{\alpha\beta} + \left.\frac{\partial P^\alpha}{\partial\mathcal{E}^\beta}\right|_{\vect{u} = 0},
  \qquad
  e\,Z^{*,\alpha\beta}_\kappa = \Omega\left.\frac{\partial P^\beta}{\partial u^\alpha_\kappa}\right|_{\vect{\mathcal{E}} = 0}
  = \left.\frac{\partial F^\alpha_\kappa}{\partial\mathcal{E}^\beta}\right|_{\vect{u} = 0} ,
  \label{eq:E.4-epsZ}
\end{equation}
where $\Omega$ is the cell volume and $\vect{u}_\kappa$ a uniform displacement of sublattice $\kappa$; Ref.~\cite[Eqs.~(94)--(96)]{baroni2001phonons} writes the same relations in Gaussian units, with $4\pi$ in front of $\partial P/\partial\mathcal{E}$. The equality of the two forms of $Z^*_\kappa$ in \cref{eq:E.4-epsZ}, as the polarization per displacement and as the force per field, holds because both are the mixed second derivative of one energy, so a code may compute whichever is cheaper. The quantities $\varepsilon_\infty$ and $Z^*_\kappa$ are all that the nonanalytic part of the dynamical matrix needs (\cref{sec:5.5}): in Heaviside--Lorentz units
\begin{equation}
  D^{\mathrm{dyn,na},\alpha\beta}_{\kappa\kappa'}(\vect{Q}\to 0) = \frac{e^2}{\Omega\sqrt{M_\kappa M_{\kappa'}}}\,
  \frac{(Z^*_\kappa\cdot\vect{Q})_\alpha\,(Z^*_{\kappa'}\cdot\vect{Q})_\beta}{\vect{Q}\cdot\varepsilon_\infty\cdot\vect{Q}} ,
  \label{eq:E.4-na}
\end{equation}
which is Ref.~\cite[Eq.~(98)]{baroni2001phonons} with $4\pi e^2_{\mathrm{G}} \to e^2$ and the mass factors of \cref{eq:A.6-Ddyn}. The same two tensors build the long-range vertex $g^{\mathrm{L}}$ of \cref{sec:6.2} and of \cref{app:E.5}.

Measured frequencies give an independent check on $Z^*$. For a cubic crystal with two atoms per cell, $Z^*_1 = -Z^*_2 = Z^*$ times the identity, and \cref{eq:E.4-na} gives the LO--TO splitting of \cref{sec:5.5},
\begin{equation}
  \omega_{\mathrm{L}}^2 - \omega_{\mathrm{T}}^2 = \frac{e^2Z^{*2}}{\Omega\,\varepsilon_\infty\,\mu_{\mathrm{red}}} ,
  \label{eq:E.4-LOTO}
\end{equation}
with $\mu_{\mathrm{red}} = M_1M_2/M_{\mathrm{cell}}$ the reduced mass of the pair. With the measured $\omega_{\mathrm{L}}$, $\omega_{\mathrm{T}}$ and $\varepsilon_\infty$ (or $\varepsilon_0$, through the Lyddane--Sachs--Teller relation of \cref{sec:5.5}), \cref{eq:E.4-LOTO} fixes $|Z^*|$. Where a rate depends on a near-cancellation, as the electron-coupled scalar in GaAs does through the vertex $Z_\kappa - Z^*_\kappa$ (\cref{sec:11.8}), this check matters more than any convergence parameter. A rate $\propto (Z_\kappa - Z^*_\kappa)^2$ has relative error $2\,\delta Z^*_\kappa/|Z_\kappa - Z^*_\kappa|$, so a one-percent error in $Z^*_\kappa$ becomes a large error in the rate when $|Z_\kappa - Z^*_\kappa|$ is a small fraction of $Z^*_\kappa$.

The Berry-phase theory of polarization (\cref{sec:5.5}) is the other route to $Z^*_\kappa$. The electronic polarization of an insulator is
\begin{equation}
  \vect{P}_{\mathrm{el}} = -\frac{e\,g_{\mathrm{deg}}}{(2\pi)^3}\sum_{v}\int_{\mathrm{BZ}}\dd^3k\,\langle u_{v\vect{k}}|\ii\nabla_{\vect{k}}|u_{v\vect{k}}\rangle ,
  \label{eq:E.4-berry}
\end{equation}
with $g_{\mathrm{deg}} = 2$ for spin, and the integrand a Berry connection of the periodic parts $u_{v\vect{k}}$~\cite[Eqs.~(89)--(90)]{marzari2012maximally}. On a grid of $\vect{k}$ the derivative is replaced by the phase of a product of overlaps $\langle u_{v\vect{k}}|u_{v'\vect{k}+\Delta\vect{k}}\rangle$ along strings of points, and this phase is gauge invariant. A Born charge is a difference of two polarizations at small displacements of opposite sign, $\pm\delta u$, and the two must lie on the same branch of $\vect{P}_{\mathrm{el}}$, which is defined only modulo the quantum $e g_{\mathrm{deg}}\vect{R}/\Omega$ for a lattice vector $\vect{R}$~\cite[Sect.~V.A.1]{marzari2012maximally}; with $\delta u$ a small fraction of a bohr the two polarizations lie on the same branch automatically. Linear response and finite differences of \cref{eq:E.4-berry} agree within numerical accuracy~\cite[Sect.~II.D.2]{baroni2001phonons}, and the Berry-phase route is the one that also works at finite fields and for ferroelectrics.

The acoustic sum rule for charges and the comparison with the measured $\varepsilon_\infty$ cost nothing to check. The sum rule, $\sum_\kappa Z^*_\kappa = 0$, holds exactly only at convergence; its violation on a finite $\vect{k}$ grid is the most direct measure of how well $\varepsilon_\infty$ and $Z^*$ are converged. The difference between the computed and the measured $\varepsilon_\infty$ comes mainly from the functional. Local and semilocal functionals put the gap too low (\cref{sec:3.4}), the energy denominators of \cref{eq:E.4-rcomm} come out too small, and the material is overscreened. The direction of the error follows from the static limit of the Kramers--Kronig relation of \cref{sec:4.3},
\begin{equation}
  \varepsilon_1(\omega = 0) - 1 = \frac{2}{\pi}\int_0^\infty\frac{\varepsilon_2(\omega')}{\omega'}\,\dd\omega' :
  \label{eq:E.4-kk}
\end{equation}
shifting the absorption spectrum $\varepsilon_2$ down in energy by a fraction of an eV, at roughly fixed weight, increases the integral. For silicon, a DFT calculation finds $\varepsilon_\infty = 12.89$~\cite{ponce2018towards}, against the measured room-temperature value of 11.7~\cite{franta2017temperature} (\cref{sec:4.4}). When $\varepsilon_\infty$ enters a rate through \cref{eq:E.4-na} or the Fr\"ohlich vertex, the measured value is usually the better input.

\subsection{Force-constant interpolation}

The Fourier interpolation of \cref{sec:5.4} is exact only if the true force constants vanish beyond the $N_Q^3$ supercell, so the grid is chosen large enough that the force constants have decayed at the edge of the supercell. A rule of thumb is to include neighbors out to two or three bond lengths and to check the interpolated frequencies against direct DFPT runs at a few $\vect{Q}$ off the grid~\cite[Sect.~II.D.3]{baroni2001phonons}. The long-range constants along the $\langle 110\rangle$ chains of bonds (\cref{fig:5-dfpt}) make silicon need larger grids than its short bonds suggest.

In a polar crystal the nonanalytic term \cref{eq:E.4-na}, in its Ewald-summed form, is subtracted from the computed $D^{\mathrm{dyn}}(\vect{Q})$ before the transform. Adding the nonanalytic term back analytically at the target $\vect{Q}$ restores the LO--TO splitting and its dependence on the direction of $\vect{Q}\to 0$.

Codes impose the acoustic sum rule after the fact, by \cref{eq:5.4-asrfix} or by a minimal projection. The size of the correction is a diagnostic: the acoustic frequencies at $\Gamma$ before it should be small compared with the lowest frequency the application needs, and a value that grows instead of shrinking with the cutoff signals an unconverged real-space grid. The correction is applied to $D^{\mathrm{dyn}}$ and to $Z^*$, never as a zeroing of $g$ at $\vect{Q}\to 0$: the acoustic vertex vanishes there by its own sum rule (\cref{sec:6.2}), and its leading term, linear in $\vect{Q}$, is physical.

A frozen-phonon calculation of the force constants out to a range $R_{\mathrm{IFC}}$ (in units of the cell size) needs supercells of $\sim R_{\mathrm{IFC}}^3$ cells, $N_{\mathrm{at}}R_{\mathrm{IFC}}^3$ atoms; with ground-state cost cubic in the number of atoms and $3N_{\mathrm{at}}$ independent displacements, the total cost scales as $3N_{\mathrm{at}}^4R_{\mathrm{IFC}}^9$. DFPT works in the primitive cell on a grid of $\sim R_{\mathrm{IFC}}^3$ wavevectors, each costing $\sim N_{\mathrm{at}}^3$ per perturbation, which totals $\sim 3N_{\mathrm{at}}^4R_{\mathrm{IFC}}^3$~\cite[Sect.~IV.B]{baroni2001phonons}, a factor $R_{\mathrm{IFC}}^6$ less.

\subsection{Displacement correlations from DFPT phonons}

The Debye--Waller exponent $2W_\kappa(\vect{q}) = \vect{q}\cdot B_\kappa\cdot\vect{q}$ and the displacement correlation $\mathcal{C}_{\kappa\kappa'}(t)$ of \cref{sec:5.6} are sums over all phonons of the zone. Inserting the mode expansion \cref{eq:A.6-u} into $B_\kappa = \langle\vect{u}_\kappa\vect{u}_\kappa\rangle$ and using $\langle B_{\vect{Q}\nu}B_{-\vect{Q}\nu}\rangle = 2n_{\vect{Q}\nu} + 1$ in the harmonic thermal state gives
\begin{equation}
  B^{\alpha\beta}_\kappa = \frac{1}{N}\sum_{\vect{Q}\nu}\frac{2n_{\vect{Q}\nu} + 1}{2M_\kappa\,\omega_{\vect{Q}\nu}}\,
  e^\alpha_{\kappa\nu}(\vect{Q})\,e^{\beta}_{\kappa\nu}(\vect{Q})^* ,
  \label{eq:E.4-B}
\end{equation}
and $\mathcal{C}_{\kappa\kappa'}(t)$ is the same sum with the phases $\ee^{\mp\ii\omega_{\vect{Q}\nu}t}$ on the emission and absorption parts and the cross products $e_{\kappa\nu}e^*_{\kappa'\nu}$. These are post-processing sums: the frequencies and polarization vectors come from the interpolated force constants on a fine grid, and the cost is negligible. The approximation is $n_A = 0$, with the thermal occupation exact for the harmonic crystal; at high temperature or for light, soft lattices the anharmonic methods of \cref{app:E.8} replace the harmonic phonons in \cref{eq:E.4-B}.

The grid sum converges slowly because of the zone center, where the acoustic frequencies vanish as $\omega = c_\nu Q$ and the summand of \cref{eq:E.4-B} diverges. The divergence is integrable in three dimensions, but a uniform grid samples it badly. How badly follows from replacing the sum over the grid cell around $\Gamma$ by an integral over a sphere of the same volume, of radius $Q_c$ with $(4\pi/3)Q_c^3 = (2\pi)^3/(N_Q^3\Omega)$. With the acoustic polarization $\vect{e}_{\kappa\nu} = (M_\kappa/M_{\mathrm{cell}})^{1/2}\hat{\vect{e}}_\nu$ of \cref{app:A.6}, an isotropic average over directions, and $\sum_\nu$ over the three acoustic branches, the contribution of that sphere to one diagonal component is
\begin{equation}
  \Delta B^{\alpha\alpha} = \frac{\Omega}{2\pi^2}\int_0^{Q_c}\!Q^2\,\dd Q\;\frac13\sum_\nu\frac{2n_{\vect{Q}\nu}+1}{2M_{\mathrm{cell}}c_\nu Q}
  \simeq
  \begin{cases}
    \dfrac{\Omega\,Q_c^2}{24\pi^2M_{\mathrm{cell}}}\displaystyle\sum_\nu\frac{1}{c_\nu}, & T = 0,\\[2ex]
    \dfrac{\Omega\,Q_c\,T}{6\pi^2M_{\mathrm{cell}}}\displaystyle\sum_\nu\frac{1}{c_\nu^2}, & T \gg c_\nu Q_c,
  \end{cases}
  \label{eq:E.4-dBgamma}
\end{equation}
where the high-temperature form uses $2n + 1 \simeq 2T/\omega$. Since $Q_c \propto 1/N_Q$, the error of a calculation that simply omits the $\Gamma$ cell falls as $N_Q^{-2}$ at zero temperature but only as $N_Q^{-1}$ at room temperature. For silicon, with $\Omega = a^3/4 = \SI{40.05}{\angstrom^3}$, $M_{\mathrm{cell}}$ two silicon masses, a representative sound speed of $\SI{6}{km.s^{-1}}$ for all three branches (the range is 3 to \SI{10}{km.s^{-1}} across materials, \cref{sec:1.3}) and an $8^3$ grid, $Q_c = \SI{0.142}{\angstrom^{-1}}$ and the room-temperature term is $\Delta B^{\alpha\alpha} \approx \SI{3.6e-4}{\angstrom^2}$, while the zero-temperature term is twenty times smaller, \SI{1.9e-5}{\angstrom^2}. (The high-temperature form applies since $c\,Q_c$ corresponds to \SI{65}{K}.) The remedy is to add \cref{eq:E.4-dBgamma}, evaluated with the computed sound speeds, as an analytic correction, or to use a grid fine enough that the residual is below the target accuracy. The same care applies to $\mathcal{C}_{\kappa\kappa'}(t)$ at long times, which are controlled by the smallest frequencies.

\subsection{Wavefunction changes beyond the density}

The vertex of \cref{eq:A.6-g} needs only $\partial_{\kappa,\vect{Q}}V_{\mathrm{KS}}$ and two Bloch states, so every DFPT run already contains it~\cite[Sect.~III.B.3]{giustino2017electron}. Codes write the potentials to disk for each irreducible $\vect{Q}$ of the coarse grid, one complex function on the real-space grid per perturbation, and a later step forms the matrix elements and interpolates them (\cref{app:E.5}). The storage is modest by modern standards: for a two-atom cell, the six perturbations per $\vect{Q}$, each a complex function on a real-space grid of $\sim 10^5$ points, take about \SI{10}{MB} per wavevector, so a few hundred irreducible wavevectors fit in a few gigabytes.

A transition amplitude needs more than the density does. The density vertex of \cref{sec:3.2}, $\rho_{cv}(\vect{k},\vect{q}) = \langle\psi_{c,\vect{k}+\vect{q}}|\ee^{\ii\vect{q}\cdot\vect{r}}|\psi_{v\vect{k}}\rangle$, changes under a frozen displacement of crystal momentum $\vect{Q}$ by
\begin{equation}
  \partial\rho_{cv} = \langle\partial\psi_{c}|\ee^{\ii\vect{q}\cdot\vect{r}}|\psi_{v}\rangle + \langle\psi_{c}|\ee^{\ii\vect{q}\cdot\vect{r}}|\partial\psi_{v}\rangle ,
  \label{eq:E.4-drho}
\end{equation}
which is the adiabatic (Herzberg--Teller) form of the phonon-assisted amplitude of \cref{sec:13.2}, valid to relative order $\omega_{\vect{Q}\nu}/\delta_{\mathrm{off}}$, with $\delta_{\mathrm{off}}$ the energy by which the intermediate states are off shell. Both terms need the complete first-order change of an individual state, while \cref{eq:5.4-sth} returns only part of $\partial\psi_{v}$ and nothing of $\partial\psi_{c}$. The full change of an occupied state is
\begin{equation}
  |\partial\psi_{v\vect{k}}\rangle = \sum_{m\notin\mathrm{deg}(v\vect{k})}|\psi_{m,\vect{k}+\vect{Q}}\rangle\,
  \frac{\langle\psi_{m,\vect{k}+\vect{Q}}|\partial_{\kappa,\vect{Q}}V_{\mathrm{KS}}|\psi_{v\vect{k}}\rangle}{\varepsilon_{v\vect{k}} - \varepsilon_{m,\vect{k}+\vect{Q}}} ,
  \label{eq:E.4-dpsi}
\end{equation}
with the states degenerate with $\psi_{v\vect{k}}$ excluded. DFPT computes the part with $m$ empty, by \cref{eq:5.4-sth}, without ever forming the sum. The part with $m$ occupied, $\op{P}^{\mathrm{occ}}_{\vect{k}+\vect{Q}}|\partial\psi_{v\vect{k}}\rangle$, is left out because it cancels in the density (\cref{sec:5.4}). In \cref{eq:E.4-drho} these terms do not cancel, because $\psi_c$ singles out one state. They are, however, cheap. The sum runs over the few occupied bands at $\vect{k}+\vect{Q}$, and the numerators are the occupied--occupied matrix elements of the potential that DFPT has already stored. The denominators $\varepsilon_{v\vect{k}} - \varepsilon_{v',\vect{k}+\vect{Q}}$ vanish wherever two valence bands at $\vect{k}$ and $\vect{k}+\vect{Q}$ cross, and near these crossings the adiabatic form fails: there the dynamical denominators $\varepsilon_{v\vect{k}} - \varepsilon_{v',\vect{k}+\vect{Q}} \pm \omega_{\vect{Q}\nu}$ of the full amplitude (\cref{sec:13.2,sec:13.6}) must be used instead. At $\vect{Q} = 0$, moreover, the degenerate subspace has to be treated in the same gauge as the states that enter $\rho_{cv}$ (\cref{app:A.5}).

Standard DFPT does not produce the change of an empty state, since empty states do not contribute to the density. The change of an empty state obeys a Sternheimer equation with the converged potential on the right,
\begin{equation}
  \bigl(H_{\mathrm{KS}} - \varepsilon_{c\vect{k}}\bigr)|\partial\psi_{c\vect{k}}\rangle = -\bigl(1 - \op{P}^{\mathrm{deg}}_{c}\bigr)\,\partial_{\kappa,\vect{Q}}V_{\mathrm{KS}}\,|\psi_{c\vect{k}}\rangle ,
  \label{eq:E.4-dpsic}
\end{equation}
with $\op{P}^{\mathrm{deg}}_{c}$ the projector on the states at $\vect{k}+\vect{Q}$ degenerate with $\psi_{c\vect{k}}$. No self-consistency is needed, because $\partial V_{\mathrm{KS}}$ is already the screened potential of the occupied manifold, and the absence of a self-consistency loop is the main saving. The operator on the left, however, is indefinite: $\varepsilon_{c\vect{k}}$ lies inside the spectrum of $H_{\mathrm{KS}}$, with all the valence states and part of the conduction bands below it, and the conjugate-gradient solvers used for \cref{eq:5.4-sth} do not converge on an indefinite operator. One method splits the space. A window $\mathcal{W}$ of bands at $\vect{k}+\vect{Q}$ contains every state below $\varepsilon_{c\vect{k}}$ and the conduction bands up to some energy above $\varepsilon_{c\vect{k}}$; the window is summed explicitly with \cref{eq:E.4-dpsi}, where the Wannier-interpolated vertices of \cref{app:E.5} make the terms cheap. The remainder of $\partial\psi_{c\vect{k}}$ obeys
\begin{equation}
  \op{P}^{\perp}_{\mathcal{W}}\bigl(H_{\mathrm{KS}} - \varepsilon_{c\vect{k}}\bigr)\op{P}^{\perp}_{\mathcal{W}}\,|\partial\psi^{\perp}_{c\vect{k}}\rangle = -\op{P}^{\perp}_{\mathcal{W}}\,\partial_{\kappa,\vect{Q}}V_{\mathrm{KS}}\,|\psi_{c\vect{k}}\rangle ,
  \label{eq:E.4-dpsiperp}
\end{equation}
with $\op{P}^{\perp}_{\mathcal{W}}$ the projector on the complement of the window. Every state left in the complement lies above $\varepsilon_{c\vect{k}}$, so the projected operator is positive definite and the standard solvers apply. The remainder contains the high-energy intermediate states, which matter for large momentum transfers, where $\ee^{\ii\vect{q}\cdot\vect{r}}$ couples to bands far above the gap. The window is a convergence parameter: the sum of the explicit and the Sternheimer parts must not depend on it.

The same equations at finite frequency give a dynamical response without empty states. For a perturbation oscillating as $\ee^{-\ii\omega t}$, first-order time-dependent perturbation theory (\cref{sec:2.4}) leads to two components of the change of each occupied state,
\begin{equation}
  \bigl(H_{\mathrm{KS}} - \varepsilon_{v\vect{k}} \mp \omega\bigr)|\partial\psi^{\pm}_{v\vect{k}}\rangle = -\bigl(1 - \op{P}^{\mathrm{occ}}\bigr)\,\partial V_{\mathrm{KS}}(\pm\omega)\,|\psi_{v\vect{k}}\rangle ,
  \label{eq:E.4-sthw}
\end{equation}
with the density change built from both components and the potential updated as in \cref{eq:5.4-dV} with an adiabatic kernel. Below the gap the operators remain positive definite on the empty manifold, because $\varepsilon_c - \varepsilon_v - \omega > 0$ for every pair, and the solution is as cheap as a static one. Above the gap a small broadening $\eta_{\mathrm{b}}$, $\omega \to \omega + \ii\eta_{\mathrm{b}}$, is added and the linear systems become complex and non-Hermitian. Frequency-dependent Sternheimer calculations of this kind give the inverse frequency moments and the drag mass of \cref{sec:12.5}, and the screened interaction in Sternheimer formulations of the methods of \cref{app:E.3}.

\subsection{Elastic constants}

The sound speeds of \cref{sec:11.1} follow from the elastic tensor $C_{ijkl}$ through the Christoffel equation of \cref{app:C.18}; for a cubic crystal along $[100]$ it reduces to $\rho_T c_{\mathrm{LA}}^2 = C_{11}$ and $\rho_T c_{\mathrm{TA}}^2 = C_{44}$ in Voigt notation. The slopes of the interpolated acoustic branches are the cheapest route to $C_{ijkl}$ and the least accurate, since the small-$Q$ part of the dispersion is precisely where the Fourier interpolation of \cref{eq:E.4-ifc} is weakest, and the slopes converge only slowly with $N_Q$. Finite strains are the most reliable route: strains of a few tenths of a percent of both signs are applied, the internal coordinates are relaxed at each strain, and the stress is fitted. The internal relaxation matters for any structure with internal degrees of freedom, and in the diamond structure it affects $C_{44}$ because the two sublattices shift against each other under shear. Since relaxation can only lower the energy at fixed strain, the relaxed constants are never larger than the clamped ones. DFPT with strain as the perturbation returns the clamped-ion tensor and the internal-strain couplings in one run, from which the relaxed tensor follows by eliminating the internal displacements~\cite[Sect.~II.E.1]{baroni2001phonons}.

A cheap check against the Voigt-order error of \cref{sec:11.1} is to compute $\rho_Tc^2$ from the tensor along one direction and compare it with the slope of the computed dispersion along the same direction.

\subsection{Convergence parameters and cost}

Five numerical parameters control a DFPT calculation. Each must be converged on the quantity of interest, not on the total energy, and with a tolerance fixed before the scan. The plane-wave cutoff of \cref{app:E.1} sets the basis; the test is the frequency at $\Gamma$ and at one zone-boundary point at two cutoffs. The $\vect{k}$ grid sets the Brillouin-zone sum in \cref{eq:5.4-dn}; for an insulator the density and the frequencies converge on modest grids, but $\varepsilon_\infty$ and $Z^*$ need noticeably finer ones, because the response to a field involves the velocity matrix elements and energy denominators of \cref{eq:E.4-rcomm} at every $\vect{k}$, and the $\vect{k}$ sum of these terms converges slowly where the direct gap is small~\cite[Sect.~II.C.2]{baroni2001phonons}. Metals need smearing of the occupations (\cref{app:E.1}) and a correspondingly finer grid. The $\vect{Q}$ grid sets the range of the interpolated force constants and is tested by direct runs off the grid. The self-consistency threshold of \cref{eq:5.4-dV} must be much tighter than in a ground-state run, because an error in the first-order density enters the force constants linearly (\cref{sec:5.4}). Where a frequency is small, as for the acoustic modes near $\Gamma$ or a soft mode, a small error in the force constant is a large relative error in the frequency, or makes it imaginary.

The lattice constant is also a parameter. Phonons are usually computed at the theoretical equilibrium of the functional. In diamond and zincblende crystals the forces vanish by symmetry at any lattice constant, but only at the equilibrium does the stress vanish, so that the second derivatives are taken at a minimum of the same energy surface. Each frequency moves with the lattice constant through its Gr\"uneisen parameter (\cref{app:E.8}); the generalized-gradient approximation, which expands lattice constants, lowers the branches with positive Gr\"uneisen parameter, and for silicon the local-density approximation already reproduces the measured dispersion closely while the generalized-gradient approximation slightly worsens the agreement~\cite[Sect.~V.A]{baroni2001phonons}.

The cost of a complete phonon calculation is the number of irreducible wavevectors on the $\vect{Q}$ grid, times $3N_{\mathrm{at}}$ perturbations (fewer, by symmetry, at high-symmetry wavevectors), times the cost of a ground-state calculation with the reduced symmetry of the perturbed crystal. The reduced symmetry raises the cost: a perturbation at a general $\vect{Q}$ leaves few symmetry operations, so the $\vect{k}$ sum runs over many more points than in the ground state. An $N_Q^3$ grid has about $N_Q^3/48$ irreducible points for a cubic crystal when $N_Q$ is large; an $8^3$ grid for silicon has 29, a count that follows from applying the 48 operations of the cubic group to the grid points and counting the distinct orbits. With six perturbations per wavevector, fewer at the high-symmetry points, and the three field perturbations at $\Gamma$, a complete silicon calculation on that grid costs at most about 180 self-consistent linear-response solutions, each of the order of one ground-state calculation at reduced symmetry. The cost grows with the fourth power of the number of atoms per cell~\cite[Sect.~IV.B]{baroni2001phonons}.

\begin{workedexample}{A DFPT calculation of silicon}\label{we:E-dfpt-si}
The settings below come from two published silicon calculations. The first used the local-density approximation with plane-wave pseudopotentials, a \SI{35}{Ry} cutoff and a $14^3$ grid for the ground-state density, and computed the dynamical matrices and the potential changes on a $6^3$ grid of wavevectors~\cite{noffsinger2012phonon}. The second used the generalized-gradient approximation at \SI{35}{Ry} (\SI{45}{Ry} for the local-density approximation), again a $6^3$ wavevector grid, and found that an $18^3$ grid of $\vect{k}$ points was needed for the Born charges to vanish~\cite{ponce2018towards}. A \SI{35}{Ry} cutoff corresponds to plane waves up to $|\vect{G}| = (2m_e\times\SI{476}{eV})^{1/2} = \SI{22}{keV}$, between five and six times the shortest reciprocal-lattice vector $|\vect{G}_{111}| = \sqrt3\,2\pi/a = \SI{3.95}{keV}$ (\cref{sec:3.1}). This cutoff is ample for phonons but, as \cref{app:E.6,app:E.7} explain, not automatically for a density vertex at large momentum transfer.

The vanishing of the Born charges is a symmetry statement that the numerics must reproduce. The two atoms of the diamond structure are exchanged by inversion through the bond center. $Z^*_\kappa$ is the derivative of a polar vector with respect to a polar vector, so it is even under inversion, and $Z^*_1 = Z^*_2$. Each atom occupies a site of tetrahedral symmetry, so each tensor is a multiple of the identity. Charge neutrality, $Z^*_1 + Z^*_2 = 0$, then forces both to vanish. On a coarse $\vect{k}$ grid the computed Born charges do not vanish, and since $Z^*$ enters \cref{eq:E.4-na}, a residual Born charge produces a spurious LO--TO splitting in a crystal that has none. The $\vect{k}$ grid is therefore fixed by the Born charge, not by the frequencies.

The dielectric constant is the second check, and here the functional limits the result. The PBE calculation of \cref{we:E-si-bands} reaches 13.04 at the measured lattice constant on an $18^3$ grid, within about one percent of the 12.89 of Ref.~\cite{ponce2018towards}. The \SI{11}{\percent} excess over the measured 11.7 is the overscreening of \cref{eq:E.4-kk} and does not shrink with any numerical parameter. The frequencies are the third check. The measured optical frequency at $\Gamma$ is \SI{15.5}{THz}~\cite{nilsson1972study}, and the local-density dispersion agrees with neutron scattering across the zone~\cite[Sect.~V.A]{baroni2001phonons}.
\end{workedexample}
